% ============================================================
% Appendix
% ============================================================

\section{Algorithm Pseudocode}
\label{app:algorithm}

\begin{algorithm}[h]
\DontPrintSemicolon
\SetKwInOut{Input}{Input}
\SetKwFor{For}{for}{do}{end}
\SetKwIF{If}{ElseIf}{Else}{if}{then}{else if}{else}{end}
\caption{HAT --- Hippocampus-Augmented Transformer Self-Improvement Loop}
\label{alg:hat_appendix}
\Input{Base model $\modeltheta$ (Cortex), Hippocampus Agent $\hippocampus$, Oracle $\oracle$, empty Neocortex $\neocortex \leftarrow \emptyset$}
\Input{SWS interval $N$, write threshold $\tau$, oracle threshold $\tau_\oracle$, selection weights $\alpha, \beta, \gamma$}
$t \leftarrow 0$\;
\For{each user interaction $(\context, \seq)$}{
  $\seqy \leftarrow \modeltheta(\context, \seq)$ \tcp*{Cortex generates response}
  $\uncertainty \leftarrow \entropy\!\big(p_\params(\seqy \mid \context, \seq)\big)$ \tcp*{Estimate uncertainty}
  Observe user feedback $\feedback$ (correction signal, if any)\;
  \tcp{--- Hippocampus Agent ---}
  $\memtrace \leftarrow \hippocampus_{\mathrm{abs}}(\context, \seq, \seqy)$ \tcp*{Abstraction}
  \If{$\uncertainty > \tau_\oracle$ \textbf{or} $\feedback = 1$}{
    $\seqy^\oracle \leftarrow \oracle(\context, \seq)$ \tcp*{Consult Oracle}
    Update $\memtrace$ with oracle response $\seqy^\oracle$\;
  }
  Compute novelty: $\novelty \leftarrow \mathrm{NoveltyScore}(\memtrace, \neocortex)$\;
  Compute selection score: $\selectscore \leftarrow \alpha\,\uncertainty + \beta\,\feedback + \gamma\,\novelty$\;
  \If{$\selectscore > \tau$}{
    $(\seq^{\mathrm{train}}, \seqy^{\mathrm{train}}) \leftarrow \hippocampus_{\mathrm{replay}}(\memtrace)$ \tcp*{Replay Builder}
    $\neocortex \leftarrow \neocortex \cup \{(\seq^{\mathrm{train}}, \seqy^{\mathrm{train}}, \selectscore)\}$ \tcp*{Write to Neocortex}
  }
  $t \leftarrow t + 1$\;
  \tcp{--- Slow-Wave Sleep ---}
  \If{$t \bmod N = 0$}{
    Sample replay batch $\replaybuf \sim \neocortex$ (priority sampling by $\selectscore$)\;
    $\params \leftarrow \params - \eta\,\nabla_\params\!\left[\losssws(\params; \replaybuf) + \lossreg(\params)\right]$ \tcp*{SWS update}
    Update Fisher information $F$ for EWC\;
  }
}
\Return Updated model $\modeltheta$\;
\end{algorithm}

\section{Hyperparameter Sensitivity}
\label{app:hyperparameters}

\begin{table}[h]
  \centering
  \caption{Hyperparameter sensitivity analysis on the NQ benchmark (Phi-3-3.8B).}
  \label{tab:hyperparameters}
  \small
  \begin{tabular}{@{}lccc@{}}
    \toprule
    \textbf{Hyperparameter} & \textbf{Range Tested} & \textbf{Best Value} & \textbf{Sensitivity} \\
    \midrule
    Write threshold $\tau$              & $\{0.1, 0.2, 0.3, 0.5, 0.7\}$ & 0.3 & Medium \\
    Oracle threshold $\tau_\oracle$     & $\{0.5, 0.6, 0.7, 0.8, 0.9\}$ & 0.7 & Low \\
    Uncertainty weight $\alpha$         & $\{0.2, 0.3, 0.4, 0.5, 0.6\}$ & 0.4 & Medium \\
    Feedback weight $\beta$             & $\{0.2, 0.3, 0.4, 0.5, 0.6\}$ & 0.4 & Low \\
    Novelty weight $\gamma$             & $\{0.1, 0.2, 0.3, 0.4\}$       & 0.2 & Low \\
    LoRA rank                           & $\{4, 8, 16, 32\}$             & 16  & Low \\
    SWS learning rate $\eta$            & $\{1, 2, 5\} \times 10^{-5}$   & $2 \times 10^{-5}$ & Medium \\
    EWC coefficient $\mu$               & $\{0.01, 0.05, 0.1, 0.5, 1.0\}$ & 0.1 & Medium \\
    SWS interval $N$                    & $\{100, 250, 500, 1000, 2000\}$ & 500 & Low \\
    \bottomrule
  \end{tabular}
\end{table}

% TODO: Add sensitivity curves (accuracy vs. each hyperparameter).

\section{Additional Experimental Results}
\label{app:additional_results}

% TODO: Add per-task breakdowns, per-cycle learning curves, additional model sizes.

\section{Prompts Used by the Hippocampus Agent}
\label{app:prompts}

% TODO: Include the exact prompts used for abstraction, replay construction, and oracle consultation.
% This is important for reproducibility.

\section{Computational Cost Analysis}
\label{app:compute}

\begin{table}[h]
  \centering
  \caption{Computational overhead of HAT components (Phi-3-3.8B on a single NVIDIA A100 GPU).}
  \label{tab:compute}
  \small
  \begin{tabular}{@{}lcc@{}}
    \toprule
    \textbf{Component} & \textbf{Time per Interaction} & \textbf{Time per SWS Cycle} \\
    \midrule
    Cortex inference             & --\,ms  & --- \\
    Hippocampus (abstraction)    & --\,ms  & --- \\
    Hippocampus (selection)      & --\,ms  & --- \\
    Oracle consultation          & --\,ms  & --- \\
    SWS training (3 epochs)      & ---     & --\,min \\
    \midrule
    \textbf{Total overhead}      & --\,ms  & --\,min \\
    \bottomrule
  \end{tabular}
\end{table}

% TODO: Fill in measured latencies.

\section{Broader Impact Statement}
\label{app:broader_impact}

HAT enables on-device language models to improve from user interactions, which raises considerations in both positive and negative directions.

\paragraph{Positive impacts.}
HAT democratises model adaptation by allowing lightweight models on personal devices to improve without requiring cloud-based retraining.
This can reduce reliance on centralised AI services and give users more control over their AI assistants.
The selective memory mechanism also implies that not all data is stored, which can be beneficial for privacy.

\paragraph{Potential risks.}
Storing user interactions, even in abstracted form, poses privacy risks.
If the Neocortex memory is compromised, sensitive information could be exposed.
Additionally, a model that learns from user corrections could amplify biases if the corrections themselves are biased.
We recommend deploying HAT with differential privacy mechanisms and content filtering to mitigate these risks.
